\section{Tensor packing}\label{sec:packing}

The closure proof uses a deterministic fact: a subspace of moderately
growing dimension captures a substantial part of only a small fraction
of an almost orthogonal tensor family.  We reproduce the argument of
\cite[Theorem~4.1]{FourfoldCompanion}, with the same quantitative
estimate.  Throughout this appendix, Hilbert spaces are real, tensor
products are Hilbert tensor products, and logarithms are natural.

\begin{theorem}[Tensor packing]\label{thm:packing}
There is a universal constant $C$ with the following property.  Let
$0<\varepsilon\le1$, and let $d\ge1$ and $L\ge0$ be integers.
Let $x_1,\ldots,x_m$ and $y_1,\ldots,y_m$ lie in the unit balls of
real Hilbert spaces $H_1,H_2$.  For each family suppose that a graph
on $\{1,\ldots,m\}$ of maximum degree at most $L$ contains every
off-diagonal pair whose absolute inner product exceeds $\eta$, where
\[
 0\le\eta\le\varepsilon^2/10^4.
\]
The two exceptional graphs may differ.  For every orthogonal projection
$P$ on $H_1\otimes H_2$ of rank at most $d$,
\begin{equation}\label{eq:packing-quantitative}
 \#\{j:\norm{P(x_j\otimes y_j)}\ge\varepsilon\}
 \le C(1+L)\varepsilon^{-2}
       \exp\!\left(C\varepsilon^{-2}\sqrt d\log(2+d)\right).
\end{equation}
For $d=d_m$ as defined in \eqref{eq:channel-dimension}, this is
$o_{\varepsilon,L}(m)$, uniformly over the Hilbert spaces, vectors,
projections, and allowed values of $\eta$.
\end{theorem}

The proof separates directions on which the testing space has large
variance and applies Gaussian comparison to the remaining directions.
The comparison is the finite Sudakov--Fernique inequality; see
\cite[Section~1, Equation~(1.2)]{Fernique1974} and
\cite{Sudakov1971}.  The following matrix estimate uses the Gaussian
integration-by-parts moment argument behind matrix Khintchine
inequalities; compare \cite[Lemma~6.1 and Proposition~7.3]{Tropp2018}.
We include its proof to retain dependence on the covariance trace,
rather than on the ambient dimensions.

\begin{lemma}[Gaussian matrix estimate]\label{lem:packing-gaussian}
Let $B_1,\ldots,B_r$ be matrices between two finite-dimensional real
Hilbert spaces, and put
\[
 H_j=\begin{pmatrix}0&B_j\\B_j^*&0\end{pmatrix},
 \qquad S=\sum_{j=1}^rH_j^2.
\]
If $\norm{S}_{\rm op}\le\sqrt d$ and $\tr S\le2d$ for some
$d\ge1$, then for independent standard normal variables $g_j$,
\begin{equation}\label{eq:packing-gaussian-bound}
 \E\left\|\sum_{j=1}^rg_jB_j\right\|_{\rm op}
       \le C d^{1/4}\sqrt{\log(2+d)}.
\end{equation}
The constant is independent of the dimensions of the two spaces.
\end{lemma}

\begin{proof}
Write $H=\sum_jg_jH_j$.  Gaussian integration by parts, followed by
differentiation of the matrix power, gives for every integer $q\ge2$
\[
 \E\tr H^{2q}
 =\sum_{j=1}^r\sum_{a=0}^{2q-2}
      \E\tr(H_jH^aH_jH^{2q-2-a}).
\]
For a fixed self-adjoint $H$ with eigenvalues $\lambda_u$ and a
self-adjoint matrix $K$, the absolute value of such a trace is at most
\[
 \sum_{u,v}|K_{uv}|^2
       |\lambda_v|^a|\lambda_u|^{2q-2-a}
 \le\tr(K^2|H|^{2q-2}).
\]
The inequality follows from weighted arithmetic--geometric mean and
the symmetry $|K_{uv}|=|K_{vu}|$.  Summing over $j$ gives
\begin{equation}\label{eq:packing-moment-recursion}
 \E\tr H^{2q}
 \le(2q-1)\norm{S}_{\rm op}\E\tr|H|^{2q-2}.
\end{equation}
Since $\E\tr H^2=\tr S$, iteration yields
\[
 \E\tr H^{2q}
 \le(2q-1)!!\,\norm{S}_{\rm op}^{q-1}\tr S.
\]
The operator norm of $H$ equals that of $\sum_jg_jB_j$.
Taking a $2q$-th root and using the hypotheses bounds its expectation by
\[
 C\sqrt q\,d^{1/4}(2\sqrt d)^{1/(2q)}.
\]
Choose an integer $q\ge2$ comparable to $\log(2+d)$ to obtain
\eqref{eq:packing-gaussian-bound}.
\end{proof}

\begin{proof}[Proof of Theorem~\ref{thm:packing}]
We first record an elementary consequence of the exceptional-graph
hypothesis.  If unit-ball vectors $v_j$ have absolute inner products
at most $\eta$ outside a graph of degree $L$, then for every nonempty
set $A$ of indices,
\begin{equation}\label{eq:packing-cluster}
 \left\|\frac1{|A|}\sum_{j\in A}v_j\right\|^2
       \le\frac{1+L}{|A|}+\eta.
\end{equation}
Indeed the expansion contains at most $(1+L)|A|$ ordered diagonal or
exceptional pairs.  A lower bound on the norm of an average therefore
bounds the number of vectors in that average.

It suffices to work in the finite spans of the $x_j$ and $y_j$.
To see this, let $Q$ project onto their tensor product and replace the
testing subspace $\operatorname{Ran}P$ by $Q(\operatorname{Ran}P)$.
For any vector $z$ in that tensor product, the norm of its projection
onto the new subspace is at least $\norm{Pz}$: in the supremum defining
$\norm{Pz}$, replace each unit testing vector by its $Q$-image, whose
norm is at most one and whose pairing with $z$ is unchanged.  The
dimension does not increase.

Identify tensors with matrices from $H_2$ to $H_1$.  Choose a
Hilbert--Schmidt orthonormal basis $A_1,\ldots,A_r$ of the testing
space, where $r\le d$.  The matrices
\[
 R_1=\sum_jA_jA_j^*,\qquad R_2=\sum_jA_j^*A_j
\]
have trace $r$.  Let $E_1,E_2$ be their spectral subspaces for
eigenvalues strictly greater than $\sqrt d$; each has dimension at most
$\sqrt d$.

Discard indices for which the projection of $x_j$ onto $E_1$ has norm
greater than $\varepsilon/10$, and do the same for $y_j,E_2$.
Cover either exceptional subspace's unit ball by at most
$(1+80/\varepsilon)^{\sqrt d}$ balls of radius $\varepsilon/40$.
If one ball contains $s$ of the discarded projections, their average
has norm at least $\varepsilon/20$.  Applying
\eqref{eq:packing-cluster} to the original vectors, with
$\eta\le\varepsilon^2/10^4$, gives
$s\le C(1+L)\varepsilon^{-2}$.  The number of discarded indices is
therefore at most
\begin{equation}\label{eq:packing-discarded}
 C(1+L)\varepsilon^{-2}
   \exp\!\bigl(\sqrt d\log(1+80/\varepsilon)\bigr).
\end{equation}

Let $Q_1,Q_2$ project onto $E_1^\perp,E_2^\perp$, and put
$B_j=Q_1A_jQ_2$.  The coefficient map
$z\mapsto(\ip{A_j}{z}_{\rm HS})_{j=1}^r$ is a contraction.
For a remaining index,
\[
 \norm{x_j\otimes y_j-Q_1x_j\otimes Q_2y_j}
 \le\norm{(1-Q_1)x_j}+\norm{(1-Q_2)y_j}
 \le\varepsilon/5.
\]
Thus each remaining index counted in
\eqref{eq:packing-quantitative} has coefficient vector
\[
 v_j=(\ip{x_j}{B_ay_j})_{a=1}^r,
 \qquad\norm{v_j}\ge\varepsilon/2.
\]
The map from $x_j\otimes y_j$ to $v_j$ remains contractive.
Outside the union of the two exceptional graphs, the tensor inner
products have absolute value at most $\eta^2$, and that union has
degree at most $2L$.  A collection of these coefficient vectors in a
ball of radius $\varepsilon/8$ centered at one of them has average
norm at least $3\varepsilon/8$.  Applying
\eqref{eq:packing-cluster} to its preimage tensors bounds its size by
$C(1+L)\varepsilon^{-2}$.  A maximal $\varepsilon/8$-separated subset
therefore has cardinality $l$ such that the number of indices still
to be counted is at most $C(1+L)\varepsilon^{-2}l$.

We now bound this separated set.  For $G=\sum_ag_aB_a$, the centered
Gaussian variables $Z_j=\ip{x_j}{Gy_j}$ satisfy
\[
 \E|Z_j-Z_{j'}|^2=\norm{v_j-v_{j'}}^2.
\]
Gaussian comparison gives, when $l\ge2$,
\begin{equation}\label{eq:packing-comparison}
 c\varepsilon\sqrt{\log l}
       \le\E\max_j Z_j\le\E\norm{G}_{\rm op}.
\end{equation}
For completeness, the comparison of finite centered Gaussian families
follows by applying Gaussian integration by parts to
$\beta^{-1}\log\sum_j e^{\beta z_j}$ along their independent Gaussian
interpolation.  Its derivative is
\[
 \frac\beta4\E\sum_{j,j'}p_jp_{j'}
 \left(\E|Z_j-Z_{j'}|^2-\E|Z'_j-Z'_{j'}|^2\right),
 \qquad p_j=\frac{e^{\beta z_j}}{\sum_a e^{\beta z_a}}.
\]
It is nonnegative when the first increment variances dominate.
Letting $\beta\to\infty$ compares the expected maxima.  For
\eqref{eq:packing-comparison}, take the $Z'_j$ to be independent
normals of variance $\varepsilon^2/128$ and use
$\E\max_{j\le l}g_j\ge c\sqrt{\log l}$.

Compression gives
\[
 \sum_jB_jB_j^*\le Q_1R_1Q_1,\qquad
 \sum_jB_j^*B_j\le Q_2R_2Q_2.
\]
These operators have norms at most $\sqrt d$ and traces at most $d$.
Lemma~\ref{lem:packing-gaussian} therefore applies.  Together with
\eqref{eq:packing-comparison} it yields
\[
 \log l\le C\varepsilon^{-2}\sqrt d\log(2+d).
\]
The resulting bound holds for $l\le1$ as well.  Combine it with
\eqref{eq:packing-discarded} and absorb
$\log(1+80/\varepsilon)\le C\varepsilon^{-2}$ to prove
\eqref{eq:packing-quantitative}.

For $d=d_m$, the exponent is
$O_\varepsilon((\log m)^{3/4}\log\log m)=o(\log m)$,
which proves the asserted $o(m)$ bound.
\end{proof}

Since projection norms are at most one, the estimate also gives
\begin{equation}\label{eq:packing-average}
 \frac1m\sum_{j=1}^m\norm{P(x_j\otimes y_j)}
 \le\varepsilon+
 \frac{C(1+L)}{m\varepsilon^2}
 \exp\!\left(C\varepsilon^{-2}\sqrt{d_m}\log(2+d_m)\right)
\end{equation}
when $\operatorname{rank}P\le d_m$.  Here $\varepsilon$ and then
$\eta$ are fixed before $m$ increases; $L$ may depend on that chosen
$\eta$.  This is exactly the order of quantifiers supplied by mixing
in Proposition~\ref{prop:channel-closure}.
